\begin{table}[htbp]\centering\small
\caption{Cronología operativa. La columna de fijación es el tiempo máximo que ICNIRP admite mirando el Sol eclipsado SIN filtro; el símbolo $\infty$ significa que la dosis no alcanza el límite, no que mirar sea prudente ni cómodo. La columna de acción supone gafas certificadas ISO 12312-2.}
\label{tab:timeline}
\begin{tabular}{llS[table-format=-2.2]S[table-format=1.3]rp{5.1cm}}
\toprule
Hora CEST & Fase & {Altura (\textdegree)} & {Obsc.} & Fijación (s) & Acción \\
\midrule
19:35:28 & C1: comienza el eclipse & 14.70 & 0.000 & 4 & Filtro puesto. Sin filtro no se mira ni se fotografía. \\
19:49:59 & parcial temprana & 12.01 & 0.157 & 7 & Filtro. Fijación sin filtro ya limitada a segundos. \\
19:59:59 & parcial media & 10.17 & 0.338 & 12 & Filtro. La caída de luz aún es imperceptible a simple vista. \\
20:14:59 & parcial profunda & 7.44 & 0.666 & 57 & Filtro. Empieza a notarse el cambio de color de la luz. \\
20:24:59 & sombras en creciente & 5.64 & 0.905 & $\infty$ & Filtro. Buscar bandas de sombra y luz de rendija. \\
20:28:59 & inminente & 4.93 & 0.995 & $\infty$ & Filtro puesto. Preparar el disparo de la corona. \\
20:29:22 & C2: comienza la totalidad & 4.85 & 1.000 & $\infty$ & RETIRAR filtro solo ahora. Anillo de diamante. \\
20:29:59 & máximo & 4.75 & 1.000 & $\infty$ & Sin filtro. Corona visible. Mirar. \\
20:30:36 & C3: termina la totalidad & 4.63 & 1.000 & $\infty$ & FILTRO PUESTO ANTES de que reaparezca la fotosfera. \\
20:34:59 & parcial de salida & 3.87 & 0.906 & $\infty$ & Filtro. El Sol sigue bajando hacia el horizonte. \\
20:59:59 & ocaso parcial & -0.22 & 0.317 & $\infty$ & El Sol se pone todavía eclipsado; filtro hasta que desaparezca. \\
\bottomrule
\end{tabular}
\end{table}
